\documentclass[journal=asbcd6,manuscript=article]{achemso}
\usepackage{amsmath,amssymb}
\usepackage{graphicx}
\usepackage{booktabs}
\usepackage{siunitx}
\usepackage{multirow}
\usepackage[hidelinks]{hyperref}

\SectionNumbersOn

\title{A Computational Framework for Designing Evolutionary-Robust Multi-Layer Kill Switches in Engineered Bacteria}

\author{OpenResearch}
\affiliation{OpenResearch, AI-Driven Scientific Discovery}

\begin{document}

\begin{abstract}
Genetically engineered microorganisms hold promise for medicine, agriculture, and bioremediation, yet reliable biocontainment remains a critical barrier to deployment. Kill switches---genetic circuits that terminate cell viability---are limited by mutational escape, wherein loss-of-function mutations confer a selective advantage that drives population-level failure. Here we present a computational framework that models escape dynamics for six architecturally distinct kill switch mechanisms, decomposes escape rates into five mutation types (point mutations, small indels, IS~element insertions, large deletions, and recombination), and predicts combined escape rates for multi-layer designs while accounting for mechanistic correlations between layers. Calibrated against published data from three independent studies (leave-one-study-out cross-validation: mean $|\log_{10}(\text{pred/obs})| = 0.04$), the framework identifies a triple-layer combination of CRISPR-based killing, overlapping gene entanglement, and synthetic auxotrophy that achieves a predicted escape rate of $6.1 \times 10^{-18}$---ten orders of magnitude below the NIH-recommended $10^{-8}$ threshold---at a fitness cost of 13\%. Sensitivity analysis reveals that overlapping gene entanglement is the most critical yet parameter-sensitive component (local sensitivity = 1.00), while synthetic auxotrophy provides the most robust protection (sensitivity = 0.17). This framework provides a quantitative design tool for engineering biocontainment systems with predictable failure rates.
\end{abstract}

\section{Introduction}

The engineering of microorganisms for therapeutic, industrial, and environmental applications has advanced rapidly, with engineered bacteria entering clinical trials as living therapeutics\cite{Brennan2022} and being developed as biosensors for environmental monitoring. However, the deployment of genetically engineered microorganisms (GEMs) in open environments raises significant biosafety concerns that must be addressed through robust biocontainment strategies\cite{Pantoja2022,Stirling2020,Lei2022}.

Kill switches are the primary biocontainment approach: genetic circuits designed to induce cell death in response to a defined signal. Over the past five years, diverse architectures have been developed, including toxin-antitoxin systems, CRISPR-Cas9-based genome degradation\cite{Rottinghaus2022}, overlapping gene entanglement\cite{Chlebek2023,Foo2025}, synthetic auxotrophy dependent on unnatural amino acids\cite{Chang2023}, and integrase-mediated differentiation circuits\cite{Williams2022}. These approaches span three orders of magnitude in reported escape frequencies, from $\sim$$10^{-5}$ to below $10^{-9}$ per cell per generation.

The fundamental challenge is \emph{mutational escape}: because kill switches impose a metabolic burden on host cells, any mutation that inactivates the switch confers a selective advantage. Over successive generations, escaped mutants are positively selected and overtake the population. Rottinghaus~et~al.\cite{Rottinghaus2022} demonstrated that the dominant escape mechanism for CRISPR-based switches is a characteristic 25~bp deletion mediated by tandem repeat recombination in the Ptet promoter, while Chlebek~et~al.\cite{Chlebek2023} showed that gene entanglement eliminates this large-deletion class but leaves small deletions and regulator mutations as residual escape routes. Williams and Murray\cite{Williams2022} showed that integrase-mediated differentiation decouples the selective advantage of escape from the burden of transgene expression.

The NIH Guidelines for Research Involving Recombinant DNA Molecules recommend an escape frequency below $10^{-8}$ for biological containment\cite{Rottinghaus2022}. To reliably achieve this threshold, the field has increasingly turned to multi-layer strategies: combining two or more mechanistically distinct kill switches such that escape requires simultaneous inactivation of all layers\cite{Pantoja2022}. The implicit assumption is that the combined escape rate equals the product of individual rates---i.e., that the layers are statistically independent.

However, this independence assumption may not hold in practice. Kill switch components that share regulatory elements, reside on the same plasmid, or are susceptible to the same mobile genetic elements (e.g., IS~elements) may exhibit correlated escape. A single large genomic deletion or IS~element insertion could simultaneously inactivate adjacent kill switch components, violating independence and dramatically increasing the true combined escape rate.

Despite the diversity of available architectures, the field currently lacks a unified computational framework that: (1)~models all major architectures under a common mathematical formalism with decomposed mutation spectra; (2)~quantitatively predicts combined escape rates for multi-layer designs while accounting for mechanistic correlations; and (3)~provides an optimization tool for selecting the best combination given constraints on fitness cost. Recent computational work by Byrom and Darlington\cite{Byrom2025} introduced host-aware modeling for circuit longevity, and Halvorsen~et~al.\cite{Halvorsen2022} analyzed drivers of kill switch stability, but neither provides a systematic combination optimization framework.

Here we address this gap. We model six kill switch architectures with mutation-spectrum-decomposed escape rates, implement stochastic serial passage simulations with 10,000 replicates per condition, and systematically evaluate all combinations of up to five layers with mechanistic correlation modeling. We identify optimal multi-layer designs and provide quantitative guidelines for biocontainment engineering.

\section{Results}

\subsection{Single-Layer Kill Switch Characterization}

We characterized six architecturally distinct kill switch mechanisms using parameters calibrated from published experimental data (Table~\ref{tab:single}, Figure~\ref{fig:arch}A). Each architecture's escape rate was decomposed into five mutation types: point mutations, small indels ($<$50~bp), IS~element insertions, large deletions ($>$50~bp), and recombination events (Figure~\ref{fig:arch}B).

\begin{table}[h]
\caption{Single-layer kill switch escape rates and dominant mutation types.}
\label{tab:single}
\centering
\begin{tabular}{lccl}
\toprule
Architecture & Rate (cell$^{-1}$ gen$^{-1}$) & $\log_{10}$ & Dominant Mutation \\
\midrule
Toxin-antitoxin & $9.0 \times 10^{-7}$ & $-6.05$ & Small indel \\
CRISPR (1~gRNA) & $2.9 \times 10^{-5}$ & $-4.54$ & Recombination \\
CRISPR (multi-gRNA) & $2.5 \times 10^{-9}$ & $-8.60$ & Point mutation \\
Overlapping gene & $1.0 \times 10^{-7}$ & $-7.00$ & Recombination \\
Synth.\ auxotrophy & $1.0 \times 10^{-9}$ & $-9.00$ & Recombination/HGT \\
Integrase diff. & $1.0 \times 10^{-6}$ & $-6.00$ & Integrase inversion \\
\bottomrule
\end{tabular}
\end{table}

The mutation spectrum decomposition reveals distinct vulnerability profiles. For single-gRNA CRISPR switches, 69\% of escape events arise from tandem-repeat recombination---specifically the 25~bp deletion in duplicated Ptet operator sequences identified by Rottinghaus~et~al.\cite{Rottinghaus2022} This dominant mechanism is nearly eliminated in the optimized multi-gRNA design through promoter mutagenesis and SOS response gene knockouts, shifting the spectrum to a more uniform distribution dominated by residual point mutations.

Overlapping gene entanglement\cite{Chlebek2023} suppresses large deletions ($<$1\% of escape events) by coupling the toxin gene to an essential host gene, but regulator mutations and small deletions within the toxin coding sequence remain viable escape routes, contributing approximately 50\% and 30\% of events, respectively.

Stochastic serial passage simulations (10,000 replicates, $N_0 = 10^7$, $K = 10^9$, 1:100 dilution) revealed the kinetics of escape takeover (Figure~\ref{fig:arch}C--D). Fitness costs were calibrated from published competitive index data: Chlebek~et~al.\cite{Chlebek2023} measured 37--51$\times$ competitive advantage for toxin-antitoxin escapers over 13~generations, yielding a per-generation fitness cost of approximately 25\% for the toxin-antitoxin system. Architecture-specific costs ranged from 3\% (synthetic auxotrophy) to 25\% (toxin-antitoxin), reflecting the different metabolic burdens imposed by each mechanism (Table~\ref{tab:fitness}).

\begin{table}[h]
\caption{Architecture-specific fitness costs calibrated from experimental data.}
\label{tab:fitness}
\centering
\begin{tabular}{lcc}
\toprule
Architecture & Fitness Cost & Calibration Source \\
\midrule
Toxin-antitoxin & 25\% & Chlebek 2023 (competitive index) \\
CRISPR (single/multi) & 5\% & Rottinghaus 2022 (GFP $\approx$ WT) \\
Overlapping gene & 12\% & Chlebek 2023 (reduced promoter) \\
Synth.\ auxotrophy & 3\% & Minimal metabolic cost \\
Integrase diff. & 15\% & Williams 2022 (T7 RNAP burden) \\
\bottomrule
\end{tabular}
\end{table}

With these calibrated costs, the median time to 50\% escape ranged from 53~generations (toxin-antitoxin) to 792~generations (synthetic auxotrophy), consistent with published experimental timescales (see Section~\ref{sec:validation}).

\subsection{Multi-Layer Combination Analysis}

We evaluated all pairwise ($n = 10$) and triple ($n = 10$) combinations of five architectures at multiple inter-layer correlation levels (Figure~\ref{fig:heatmap}). At the reference correlation $\rho = 0.01$, six of ten pairwise combinations met the NIH threshold of $10^{-8}$, while all ten triple combinations passed.

The most effective pairwise combination was CRISPR multi-gRNA + synthetic auxotrophy ($2.5 \times 10^{-11}$), exploiting the orthogonality of genome-targeting and metabolic-dependence mechanisms. The best triple combination---CRISPR multi-gRNA + overlapping gene + auxotrophy---achieved $6.06 \times 10^{-18}$, ten orders of magnitude below the NIH threshold.

Notably, combinations involving integrase differentiation circuits at $\rho = 0.01$ were marginal: at $1.0 \times 10^{-8}$, they exactly match the threshold, reflecting the relatively high individual escape rate ($10^{-6}$) of integrase-based mechanisms.

\subsection{Inter-Layer Correlation Is the Critical Design Parameter}

The inter-layer correlation $\rho$---quantifying the probability that escape from one layer mechanistically facilitates escape from another---was the most influential parameter in our analysis (Figure~\ref{fig:corr}A).

For the pairwise combination of CRISPR multi-gRNA and overlapping gene, the combined escape rate varies from $2.6 \times 10^{-16}$ (independent, $\rho = 0$) to $1.0 \times 10^{-7}$ (fully correlated, $\rho = 1$)---a span of nine orders of magnitude. The critical correlation at which this combination loses NIH compliance is $\rho_\text{crit} = 0.102$.

In contrast, the CRISPR multi-gRNA + auxotrophy pair and the recommended triple-layer design maintain NIH compliance across the entire correlation range tested, including $\rho = 1.0$ (escape rate $= 6.06 \times 10^{-16}$ for the triple). This robustness arises from the extremely low individual escape rate of synthetic auxotrophy ($10^{-9}$), which dominates the combined rate even at maximum correlation.

A two-dimensional phase diagram (Figure~\ref{fig:corr}B) mapping the combined escape rate as a function of both $\rho$ and the overlapping gene escape rate factor reveals a broad safe operating region where the NIH threshold is maintained. This visualization provides a practical tool for assessing the impact of design uncertainties.

\subsection{Optimal Design Search}

An exhaustive search over all 31 valid combinations of up to five layers (maximum fitness cost 25\%) identified 24 designs meeting the NIH threshold at $\rho = 0.01$ (Figure~\ref{fig:pareto}). Five designs formed the Pareto front---no other design achieved both a lower escape rate and lower fitness cost.

The practical recommendation is the three-layer design of CRISPR multi-gRNA + overlapping gene + auxotrophy: at $6.06 \times 10^{-18}$ escape rate and 20\% total fitness cost, it provides an enormous safety margin. Adding a fourth layer (integrase differentiation) improves the rate to $9.62 \times 10^{-24}$ but increases the cost to 23\%, with marginal practical benefit given that containment is already assured.

\subsection{Framework Validation}

We performed leave-one-study-out cross-validation, calibrating on data from Rottinghaus~et~al.\cite{Rottinghaus2022} and predicting Chlebek~et~al.\cite{Chlebek2023} data, and vice versa (Figure~\ref{fig:valid}A). All five escape rate predictions fell within one order of magnitude of experimental measurements, with a mean absolute log-error of 0.04.

\label{sec:validation}

For time-to-escape predictions, calibration of architecture-specific fitness costs from competitive index data (Table~\ref{tab:fitness}) substantially improved accuracy. The model predicted 40~generations for the toxin-antitoxin system to reach 10\% escape, compared to 30~generations observed by Chlebek~et~al.\cite{Chlebek2023} (ratio 1.3$\times$). For the entangled overlapping gene system under auxotrophic selection, the model predicted 106~generations, consistent with the experimental observation that 7 of 10 lineages had not reached 10\% escape by 132~generations. A residual 1.3$\times$ overestimate for the toxin-antitoxin system likely reflects differences in effective population size between simulated (constant $10^7$ post-dilution) and experimental (variable) conditions.

\subsection{Ablation and Sensitivity Analysis}

Ablation analysis (Figure~\ref{fig:ablation}A) confirmed that each layer in the recommended triple design provides essential, non-redundant protection. Removing any single layer increased the combined escape rate by $4.1 \times 10^6$ to $1.7 \times 10^8$ fold. The overlapping gene and auxotrophy layers contributed nearly equally ($\sim$$10^8$-fold each), while the CRISPR multi-gRNA layer contributed $\sim$$10^6$-fold.

Local sensitivity analysis (Figure~\ref{fig:ablation}B) revealed a heterogeneous vulnerability profile. The overlapping gene layer had the highest sensitivity coefficient ($S = 1.00$): a 10-fold increase in its individual escape rate translates directly to a 10-fold increase in the combined rate. The CRISPR multi-gRNA layer showed moderate sensitivity ($S = 0.84$). Synthetic auxotrophy was the most robust component ($S = 0.17$), meaning the combined rate is relatively insensitive to uncertainty in the auxotrophy escape rate.

\subsection{Robustness Across Conditions}

To assess the generality of our conclusions, we performed systematic robustness analyses across population sizes, growth environments, deployment durations, and correlation model assumptions (Supplementary Figures~S2--S4).

\paragraph{Population size.} Varying the effective population from $10^5$ to $10^9$ cells had minimal effect on escape dynamics: the toxin-antitoxin system's time to 10\% escape varied only 1.15-fold ($40$--$46$~generations), while the triple design showed no escape events at any population size within the 1,000-generation simulation window.

\paragraph{Growth environment.} Under five environments spanning rich laboratory media to nutrient-poor soil conditions, the toxin-antitoxin system's time to 50\% escape ranged from 53 to 880~generations (17-fold variation), while the CRISPR multi-gRNA system survived beyond 1,000~generations in low-growth environments. The triple design maintained zero escape across all environments.

\paragraph{Long-term deployment.} Extending simulations to 5,000 generations (5,000 replicates) revealed a clear hierarchy (Supplementary Figure~S2). The single CRISPR multi-gRNA system showed 100\% population-level escape by 1,000~generations. The double CRISPR + auxotrophy combination reached 55\% escape probability by 5,000~generations. The triple design showed \emph{zero escape events in any of 5,000 replicates through 5,000~generations}, confirming indefinite containment within the simulation horizon.

\paragraph{Correlation model sensitivity.} Testing four alternative formulations for the combination model (Supplementary Figure~S4), we found that the triple design maintained NIH compliance ($<$$10^{-8}$) under two of four models \emph{at all tested correlation values}, including $\rho = 1$. Under the most pessimistic model (maximum of independent and worst-pair-correlated), NIH compliance was lost at $\rho = 0.107$. This demonstrates that the triple design's superiority is robust to the specific mathematical formulation of inter-layer correlation.

\section{Discussion}

We have developed a computational framework for the rational design of multi-layer biocontainment systems that addresses three key limitations of the current approach: the lack of systematic cross-architecture comparison, the naive independence assumption for combination strategies, and the absence of quantitative design optimization.

\paragraph{Key Design Principle: Mechanistic Orthogonality.}
Our central finding is that effective multi-layer biocontainment requires \emph{mechanistic orthogonality}: the combined layers must exploit fundamentally different biological mechanisms such that no single mutational event can simultaneously compromise multiple layers. The recommended triple combination achieves this through three orthogonal channels: (1)~CRISPR-based genome degradation requiring multiple simultaneous promoter mutations; (2)~gene entanglement constraining the mutational landscape through physical coupling of toxic and essential genes; and (3)~metabolic auxotrophy requiring gain-of-function evolution of an entirely new biosynthetic capability.

\paragraph{The Correlation Problem.}
The critical importance of inter-layer correlation ($\rho$) represents a key insight. Published studies implicitly assume independence when arguing that combined rates equal products of individual rates\cite{Rottinghaus2022,Chang2023}. Our analysis demonstrates that at realistic correlation values ($\rho = 0.01$--$0.1$), pairwise combinations are often marginal or inadequate, while triple-layer designs maintain robust safety margins. Practically, this means that components should be placed on separate genomic loci (not the same plasmid), should use distinct promoter systems, and should not share regulatory elements.

\paragraph{Comparison to Prior Work.}
Williams and Murray\cite{Williams2022} developed ODE and stochastic models for differentiation circuits but focused on a single architecture. Byrom and Darlington\cite{Byrom2025} introduced host-aware modeling but did not optimize combinations. Our framework is the first to provide mutation-spectrum-resolved modeling across architectures with correlation-aware combination optimization.

\paragraph{Limitations.}
Several limitations warrant discussion. First, escape rates are treated as fixed parameters calibrated from specific experimental systems (\textit{E.~coli} Nissle, \textit{P.~protegens}); rates may differ in other host strains or environmental contexts, as demonstrated by the 17-fold variation in escape dynamics across growth environments (Supplementary Figure~S3). Second, while fitness cost calibration from competitive index data reduced the time-to-escape prediction error from 6--8$\times$ to 1.3$\times$ for the toxin-antitoxin system, a residual overestimate remains, likely reflecting simplified population dynamics in our serial passage model. Third, we do not model horizontal gene transfer, which represents an additional containment failure mode typically addressed through orthogonal strategies such as Cas9-mediated degradation of mobile genetic elements\cite{Hayashi2024}. Fourth, the correlation parameter $\rho$ is a scalar approximation; in practice, mechanistic correlations depend on genomic context (distance between loci, shared IS~element targets, plasmid co-location) and should be characterized experimentally for specific designs.

\paragraph{Future Directions.}
Experimental validation of the recommended triple-layer combination is the critical next step. Integration with host-aware modeling\cite{Byrom2025} and realistic mutation spectrum simulation (tracking individual IS~element insertions, point mutations, and deletions) would further improve predictive accuracy.

\section{Methods}

\subsection{Mutation Spectrum Decomposition}
Each kill switch architecture was modeled as a single layer characterized by a mutation spectrum decomposing the total escape rate into five mechanistic classes: point mutations, small indels ($<$50~bp), IS~element insertions, large deletions ($>$50~bp), and recombination events. Individual rates were calibrated from published experimental data\cite{Rottinghaus2022,Chlebek2023,Williams2022} using the reported escape frequencies, mutation characterization by sequencing, and mechanistic analysis of escape isolates.

\subsection{Multi-Layer Combination Model}
For $n$ layers with individual escape rates $\mu_1, \ldots, \mu_n$, the combined escape rate under pairwise correlation is:
\begin{equation}
P_\text{combined} = \prod_{i=1}^n \mu_i + \sum_{i < j} \rho_{ij} \cdot \max(\mu_i, \mu_j) \cdot \prod_{k \neq i,j} \mu_k
\label{eq:combination}
\end{equation}
where $\rho_{ij}$ is the effective pairwise correlation between layers $i$ and $j$, computed as the sum of a base genomic correlation and mechanistic contributions from shared mutation vulnerabilities (IS~element targets, large deletion susceptibility).

\subsection{Stochastic Serial Passage Simulation}
Population dynamics were simulated as a serial passage process. Each passage consists of: (1)~per-generation mutation (Poisson-sampled new escape mutants from the wild-type pool); (2)~differential growth over $\sim$\,$\log_2(K/N_\text{diluted})$ generations, with wild-type cells growing at rate $2(1-c)$ per generation and escape mutants at rate 2 per generation, where $c$ is the fitness cost; (3)~carrying capacity enforcement at $K = 10^9$; and (4)~bottleneck sampling via binomial dilution to $K/100 = 10^7$ cells. Simulations used 10,000 replicates per condition with a fixed random seed (42) for reproducibility.

\subsection{Convergence and Statistical Analysis}
Convergence was verified by computing the standard error of the median time-to-escape at replicate counts from 50 to 10,000 (SE~$= 0.02$~generations at $n = 10{,}000$; Supplementary Figure~S1). Bootstrap confidence intervals (5,000 resamples) were computed for all reported statistics. Cross-validation used a leave-one-study-out protocol, calibrating on one study's data and predicting the other's.

\subsection{Code Availability}
All source code, data, and analysis scripts are available at \url{https://github.com/XiangJinyu/kill-switch-design}. Simulations were performed in Python~3.12 using NumPy~1.26 and SciPy~1.12. All results are fully reproducible from the provided scripts.

\bibliography{references}

\end{document}
